\largerpage
\sscp{\tsc{\ili{Seti}-Benggoi}}{10-03-04}
\tsc{Seti} is a boundary cluster of languages the affiliation
of which is not immediately clear due to influence from \tsc{East
Seram} languages with which it is in close contact.
In \citet[20, 37]{Collins1983},
\ili{Seti} is listed as a branch under Nunusaku,
though Collins is explicit that this placement is tentative.
In \citet{Collins1986}, \tsc{Seti} was removed from Nunusaku
and placed under his “East Seram” node.

Evidence in \citet[117--121]{LoskiLoski1989} indicates that
\tsc{Seti} is best thought of as a chain or cluster of
related languages and dialects, rather than a single language.
The percentages of lexical similarity calculated by \citet[119]{LoskiLoski1989}
are given in \trf{10-11} along with the villages surveyed in
italics and the names of the varieties identified by \citet{LoskiLoski1989} in normal typeface.
Note that their “\ili{Seti}” contains two internal groupings; \ili{Seti},
Ulahaham, and Dihil village, as well as Elnusa and Adabai villages.
With the exception of \ili{Seti} village,
all the villages that comprise these two groupings are on the
south coast of Seram,
along the northeastern coast of the \ili{Teluti} bay.

%\begin{table}
	%\centering\tcp{Lexical similarity among varieties of \tsc{Seti}}{10-11}
	%\st{6pt}\begin{tabular}{lllllllll}\lsptoprule
%\mc{3}{l}{\it{Kobi}}									&&&&&&	\ili{Kobi}	\\	\hhline{~--------}
%78	&	\mc{4}{|l}{\it{Seti}{\cgr}}&{\cgr}&{\cgr}&{\cgr}&\mc{1}{l|}{}\\
%71	&	\mc{1}{|l}{80	{\cgr}}&	\mc{3}{l}{\it{Ulahaham}}									&&&&\mc{1}{l|}{}\\
%75	&	\mc{1}{|l}{83	{\cgr}}&	83	&	\mc{3}{l}{\it{Dihil}{\cgr}}{\cgr}&{\cgr}&{\cgr}	&	\mc{1}{l|}{Seti}	\\
%69	&	\mc{1}{|l}{77	{\cgr}}&	79	&	85	{\cgr}&	\mc{3}{l}{\it{Elnusa}}							&&\mc{1}{l|}{}\\
%70	&	\mc{1}{|l}{74	{\cgr}}&	75	&	82	{\cgr}&	81	&	\mc{3}{l}{\it{Adabai}{\cgr}} 	&\mc{1}{l|}{}\\	\hhline{~--------}
%70	&	74	{\cgr}&	67	&	69	{\cgr}&	66	&	70	{\cgr}&	\mc{2}{l}{\it{Wahakaim}}			&	\ili{Wahakaim}	\\
		%\lspbottomrule
	%\end{tabular}
%\end{table}

\begin{table}
	\centering\tcp{Lexical similarity among varieties of \tsc{Seti}}{10-11}
	\st{6pt}\begin{tabular}{lllllllll}\lsptoprule
\mc{3}{l}{\it{Kobi}}									&&&&&&	\ili{Kobi}	\\	\hhline{~--------}
78	&	\mc{4}{|l}{\it{Seti}}&&&&\mc{1}{l|}{}\\
71	&	\mc{1}{|l}{80 {\cgr}}&	\mc{3}{l}{\it{Ulahaham}}									&&&&\mc{1}{l|}{}\\
75	&	\mc{1}{|l}{83 {\cgr}}&	83	{\cgr}&	\mc{3}{l}{\it{Dihil}}&&	&	\mc{1}{l|}{Seti}	\\
69	&	\mc{1}{|l}{77}&	79	&	85 	{\cgr}&	\mc{3}{l}{\it{Elnusa}}							&&\mc{1}{l|}{}\\
70	&	\mc{1}{|l}{74}&	75	&	82  {\cgr}&	81	 {\cgr}&	\mc{3}{l}{\it{Adabai}} 	&\mc{1}{l|}{}\\	\hhline{~--------}
70	&	74&	67	&	69&	66	&	70&	\mc{2}{l}{\it{Wahakaim}}			&	\ili{Wahakaim}	\\
		\lspbottomrule
	\end{tabular}
\end{table}

A number of other sources provide data from languages of eastern
Seram that appear to belong to the \tsc{Seti} cluster.
These sources are given in \trf{10-12},
below along with the label the source gives to the language they describe.
Although data in all of these sources differs to varying degrees,
it also has a level of similarity that indicates that it is probably
from a single language chain/cluster.
Of all these sources,
the “\ili{Liambata}” data in \citet{Tauern1931} is the most divergent,
and this may represent \ili{Salas} [sgu] or a variety of \ili{Benggoi} [bgy].\footnote{
		The data in \citet[128--129]{Tauern1931} is entitled: “\ili{Liambata} dialect
		of the Alifuru people of \ili{east} Seram.
		\ili{Bonfia} language”
		%(“\ili{Liambata}-Dialekt der Alfuren von Ostseran. \ili{Bonfia}-Sprache.”).
		Although in other sources “\ili{Bonfia}” is used
		to refer to \ili{Masiwang},
		the data in Tauern is definitely not \ili{Masiwang}.} 

\begin{table}
	\centering\tcp{Sources with data from \tsc{Seti}}{10-12}
	\st{3pt}\begin{tabularx}{\linewidth}{ll>{\raggedright\arraybackslash}X}\lsptoprule
Source														&	Label/Variety			&	Notes	\\	\midrule
\citet{Boot1893}									&	\ili{Kobi}							&	source of \citeauthor{Stresemann1927}'s \ili{Kobi} \\	
\citet{Wallace1869}								&	Ahtiago						&	source of \citeauthor{Stresemann1927}'s \ili{Liambata}	\\	
\citet{Stresemann1927}						&	\ili{Liambata}					&		\\	
\citet[122--128]{Tauern1931}			&	\mc{2}{l}{Uhei-Kaẖlakim}		\\	
\citet[128--139]{Tauern1931}			&	\ili{Liambata}					&	different from \citeauthor{Stresemann1927}'s \ili{Liambata}\\
\citet[177--180]{Stokhof1981a}		&	\ili{Maneo}							&		\\	
\citet[181--184]{Stokhof1981a}		&	Teula							&		\\	
\citet{Collins1986}								&	\ili{Seti} {\&} Adabai	&		\\	
\citet{McCollum-Grimes1990} 	&	Lapela-Kedalil		&		\\	
		\lspbottomrule
	\end{tabularx}
\end{table}

From west to \ili{east} along the northeastern coast of the \ili{Teluti}
bay villages where \tsc{Seti} is spoken lie in the following order
(other villages also occur in this area): Ulahaham,
Dihil, Lapela, Elnusa, \ili{Atiahu}, and Adabai.
Note that \ili{Atiahu} village lies between Elnusa
and Adabai, which together form a single grouping in \citet{LoskiLoski1989}.
\ili{Atiahu} village is presumably where the “Ahtiago” data in Wallace
1869 was collected,\footnote{
\citet{Wallace1869} also has a word list
from “Ahtiago and Tobo” which is \ili{Bobot} [bty] (\srf{11-02}).} 
which in turn appears to be the source of the “\ili{Liambata}” data in \citet{Stresemann1927}.
It is thus possible that this \ili{Atiahu}/\ili{Liambata} variety of \tsc{Seti}
belongs with Elnusa and Adabai villages as a single variety/dialect of \tsc{Seti}.

\ili{Benggoi} is another language/dialect cluster neighbouring \tsc{Seti},
and the two clusters have a number of similarities.
According to \citet[121]{LoskiLoski1989} “The \ili{Benggoi} language
is spoken by 350 people in three villages in the Werinama
and Bula districts (\ili{Benggoi}, Balakeo, and Lesa).
It is most closely related to \tsc{Seti} (59{\%}).
Although each of the dialects of \ili{Benggoi} share only 70{\%} lexical
similarity we have decided to group these as dialects of the
same language because the speakers we interviewed felt that the
language in these three villages were similar.” Our \ili{Benggoi} data
comes from \citet{ErniatiRumalean2023} and represents the variety spoke in Bula.

Unless otherwise marked, \tsc{Seti} data throughout this book is drawn
primarily from \citet{McCollum-Grimes1990},
which is from the Lapela-Kedalil dialect from the southern coast
of Seram known by its speakers as Liana.
This data is different to the \ili{Atiahu}/\ili{Liambata} data in \citet{Wallace1869}
and  \citet{Stresemann1927}.
McCollum also indicates that Ulahaham is a different dialect.
Liana thus appears to be distinct from the variety of \tsc{Seti}
spoken to its west in Ulahaham
and Dihil villages, as well as \ili{Atiahu}/\ili{Liambata} (Elnusa-Adabai?) to its \ili{east}.
%When this data needs to be distinguished from other varieties
%of \tsc{Seti} it is labelled “\ili{Liana-Seti}”.
Data in \citet{Collins1986} has some differences
and is marked as (C) throughout.

\tsc{Seti} has metathesis of the high vowel /i/,
as seen in *duRi > \it{luil-e} ‘bone’,
*tasik ‘sea, saltwater’ > \it{tais-a} ‘salt’,
and **mantəR > **mari > \it{mair-a} ‘cuscus’.
Similar metathesis of high vowels occurs in \tsc{East Seram} languages (\srf{11-02}).
In addition to having contact effects with \tsc{East Seram} languages
to its \ili{east}, \tsc{Seti} and \ili{Benggoi} have a few lexically specific innovations
which are shared with \ili{Manusela} to the west.
Examples are given in \qf{10-58}.
%\ili{Liambata} data in \qf{10-58} is from \citet{Stresemann1927}.

%\begin{table}\tcp{Lexically specific \tsc{\ili{Seti}-Benggoi} and \tsc{Manusela} innovations}{10-00}
%%\begin{adjustbox}{width=\textwidth}
	%\begin{threeparttable}\st{2pt}
	%\begin{tabular}{llllllll}\lsptoprule
%local gl. &hand	&	moon	&	rain	&	behind, back	&	hear	&	tongue	&	smoke\\
%PMP	&*qalima\su{†}	&	(*bulan)	&	(*quzan)	&	(*ma-udəhi)	&	(*dəŋəR)	&	(**maya)\su{‡}	&	(**ma-qasu)\\	\midrule
%\ili{Liana-Seti}&ima-k	&	umal-a	&	ro-a	&	kalu-k	&	au-bulai	&	lela-k	&	foina\\
%\ili{Benggoi}		&ima-m	&	umel-am	&	do-am	&		&	ya-huley	&	lela-m	&	hoin-am\\
%Mansusela	&ima-ni	&	umal-a	&	ro-a	&	kalu-nia	&	hulai	&	lela-ni	&	hoina\\
		%\lspbottomrule
	%\end{tabular}
		%\begin{tablenotes}\tnsize
			%\item[†] {Irregular *l > {\0}.}
			%\item[‡] {Replacement or PWMP *zəlay.}
		%\end{tablenotes}
	%\end{threeparttable}
%%\end{adjustbox}
%\end{table}

\noindent{\wg\st{4pt}\begin{tabularx}
{\linewidth}{>{\hspace{-\tabcolsep}}l<{\hspace{-\tabcolsep}}
>{\hspace{-\tabcolsep}\enspace}lllll
>{\raggedright\arraybackslash}X}
\lsptoprule
%\tbx{10-58}	&	PMP	&	\ili{Liana-Seti}	&	\ili{Benggoi}	&	\ili{Manusela}	&	gloss	&	innovation	\\	\midrule
%\endfirsthead
\tbx{10-58}	&	PMP	&	\ili{Liana-Seti}	&	\ili{Benggoi}	&	\ili{Manusela}	&	gloss	&		\\	\midrule
%\endhead
	&	*qalima	&	ima-k	&	ima-m	&	ima-ni	&	hand	&	irr. *l > {\0}	\\
	&	(*bulan)	&	umal-a	&	umel-am	&	umal-a	&	moon	&		\hp{\,}\\
	&	(*quzan)	&	ro-a	&	do-am	&	ro-a	&	rain	&	\hp{\,}	\\
	&	(*ma-udəhi)	&	kalu-k	&		&	kalu-nia	&	\mc{2}{l}{behind, back}	\\
	&	(*dəŋəR)	&	au-bulai	&	ya-huley	&	hulai	&	hear	&	\hp{\,}	\\
	&	(**maya)	&	lela-k	&	lela-m	&	lela-ni	&	tongue	&	PWMP **zəlaq?	\\
	&	(**ma-qasu)	&	foina	&	hoin-am	&	hoina	&	smoke	&	\hp{\,}	\\
\lspbottomrule
\end{tabularx}\mbox{}\\}

\tsc{Seti} and \ili{Benggoi} share *ŋ/*n > \it{n} in most words with \tsc{\ili{Ambon}-Seram}.
Examples are given in \qf{10-59} below.
There are two exceptions in our data which are probably attributable
to contact with neighbouring \ili{Bobot} and/or \ili{Masiwang}.
These are the reflexes of *naŋuy ‘swim’ in \tsc{Seti}
and *saŋa ‘forked branch’ in \ili{Benggoi}.

\noindent{\wg\st{2pt}\begin{tabularx}
{\linewidth}{>{\hspace{-\tabcolsep}}l<{\hspace{-\tabcolsep}}
>{\hspace{-\tabcolsep}\enspace}lll
>{\raggedright\arraybackslash}X}
\lsptoprule
%\tbx{10-59}	&	PMP	&	Liana-S.	&	\ili{Benggoi}	&	gloss	\\	\midrule
%\endfirsthead
\tbx{10-59}	&	PMP	&	Liana-S.	&	\ili{Benggoi}	&	gloss	\\	\midrule
%\endhead
	&	*\tbr{ŋ}ajan	&	\tbr{n}ala-k	&	\tbr{n}ala-m	&	name	\\
	&	*\tbr{ŋ}isi	&	\tbr{n}isi-k	&	\tbr{n}isi-r	&	tooth	\\
	&	*\tbr{ŋ}ijuŋ	&	\tbr{n}inu-k	&	\tbr{n}ilu-m	&	nose	\\
	&	**\tbr{ŋ}əli{\footnotemark} 	&	\tbr{n}el-e	&		&	pig tusk	\\
	&	**\tbr{ŋ}aRək 	&	\tbr{n}aʔa	&	\tcg{(tahun)}	&	year; (PAMS, cf. \ili{Buru} \it{ŋahe-t}, \ili{Asilulu} \it{nale}, E. \ili{Kola} \it{nahak,} SW. \ili{Tarangan} \it{narak})	\\
	&	*ləsu\tbr{ŋ}	&	lus\tbr{n}-a	&	kanlus\tbr{n}-am	&	mortar (< **losun-a)	\\
	&	*hasa\tbr{ŋ}	&	yasa\tbr{n}-a	&		&	fish gill	\\
	&	*ta\tbr{ŋ}is	&	tatá\tbr{n} (C), &	\tcg{(sela)}	&	weep	\\	\hhline{~}
	&								&	au-ra\tbr{n}	&										&	I weep (onset C alternation \srf{07-04-03})	\\
	&	*tali\tbr{ŋ}a	&	ti\tbr{n}a-k	&	ti\tbr{n}a-m	&	ear	\\
	&	*sa\tbr{ŋ}a	&	sa\tbr{nn}a	&	sok{\tl}so\tbr{k}-am	&	fork, branch (cf. \ili{Masiwang} \it{sak{\tl}sa\tbr{k}an})	\\
	&	*na\tbr{ŋ}uy	&	au-na\tbr{k}	&	\tcg{(etinem)}	&	swim (cf. \ili{Bobot} \it{na\tbr{k}u̥)}; \ili{Kobi} \ili{Seti} ‹ienè\tbr{n}› \citep[1187]{Boot1893}	\\
\lspbottomrule
\end{tabularx}\mbox{}\\}
\footnotetext{For *ŋəli ‘canine tooth’ compare \ili{Kamarian} \it{neri}
	‘tusks of pig’ \citep[315]{vanEkris1864},
	\ili{Hitu} \it{neli} ‘tusk (of a pig)’ \citep[105]{vanHoevell1877}.
	For initial *ŋ compare cognates in Timor: \ili{Helong} \it{ŋilin},
	\ili{Meto} \it{niis} \it{koni-f} ‘canine tooth’,
	Proto-Rote-\ili{Meto} *ŋoli \citep[297]{Edwards2021}.}

For the most part,
*d/*z/*l/*j > \it{l} in \tsc{Seti} and \ili{Benggoi}.
Word finally, *d/*z/*l/*j > \it{l,} {\0} with one possible case
of final *j > \it{n} in \ili{Benggoi} (*uləj > \it{ulan-em} ‘worm’).
The reflexes of *d,
*z, and *j are most straightforward with the reflex usually being \it{l}.
Examples are given in \qf{10-60}--\qf{10-62} below.
\\

\noindent{\wg\st{3pt}\begin{tabularx}
{\linewidth}{>{\hspace{-\tabcolsep}}l<{\hspace{-\tabcolsep}}
>{\hspace{-\tabcolsep}\enspace}l<{\hspace{-\tabcolsep}\,}>{\hspace{-\tabcolsep}\,}lll
>{\raggedright\arraybackslash}X}
\lsptoprule
	%&		&	PMP	&	Liana-S.	&	\ili{Benggoi}	&	gloss	\\	\midrule
%\endfirsthead	
&		&	PMP	&	Liana-S.	&	\ili{Benggoi}	&	gloss	\\	\midrule
%\endhead	\midrule
\tbx{10-60}	&	*d	&	*\tbr{d}aRaq	&	\tbr{l}aha {\tl} \tbr{l}ahe	&	\tbr{l}ahi-m	&	blood	\\
	&		&	*\tbr{d}uRi	&	\tbr{l}uil-e	&	\tbr{l}uil-am	&	bone	\\
	&		&	*\tbr{d}uha	&	en-\tbr{l}ua	&	en-\tbr{l}o	&	two	\\
	&		&	*\tbr{d}ahun	&	ai \tbr{l}ann-e	&	ai \tbr{l}an-am	&	leaf	\\
	&		&	*\tbr{d}akəp	&	au-\tbr{l}ak	&		&	I catch, seize	\\
	&		&	*\tbr{d}aləm	&	\tbr{l}ali-k, fal\tbr{l}al	&	\tcg{(dalam)}	&	inside; \ili{Benggoi} from \ili{Malay}\\
	&		&	*kə\tbr{d}əŋ	&	au k-o\tbr{l}o	&		&	stand \it{(k)olo} (C)	\\
	&		&	*tuqə\tbr{d}	&	tua-ne	&		&	tree stump	\\	\midrule
\tbx{10-61}	&	*z	&	*\tbr{z}alan	&	\tbr{l}alam (C), \tbr{l}ala	&	\tbr{l}alean	&	path	\\
	&		&	**\tbr{z}əlaq	&	\tbr{l}ela-k	&	\tbr{l}ela-m	&	tongue (PWMP)	\\
	&		&	*qa\tbr{z}ay	&	a\tbr{l}a-k	&	a\tbr{l}a-m	&	jaw, chin	\\
	&		&	*haRə\tbr{z}an	&	o\tbr{l}am (C)	&	ahe\tbr{l}a-m	&	stairway	\\	\midrule
\tbx{10-62}	&	*j	&	*pa\tbr{j}ay	&	ha\tbr{l}a	&	ha\tbr{l}e-m	&	rice	\\
	&		&	*ŋa\tbr{j}an	&	na\tbr{l}a-k, na\tbr{l}a-m	&	na\tbr{l}a-m, na\tbr{n}e-m	&	name	\\
	&		&	*pi\tbr{j}a	&	en-hi\tbr{l}a	&	in-hi\tbr{l}	&	how many?	\\
	&		&	*ma\tbr{j}a	&	ma\tbr{l}afu-a	&	me\tbr{l}a	&	Liana-S. `dry up', \ili{Benggoi} `dry'	\\
	&		&	*qapə\tbr{j}u	&	hu\tbr{l}ai-na	&		&	gall, bile (unexplained vowels)	\\
	&		&	*ŋi\tbr{j}uŋ	&	ni\tbr{n}u-k	&	ni\tbr{l}u-k	&	nose (*j > **l > \it{n} due to progressive assimilation)	\\
	&		&	*waRə\tbr{j}	&	ai wehe\tbr{l}-a	&	iwe\tbr{l}-am	&	vine, rope	\\
	&		&	*qalə\tbr{j}aw	&	lea	&	\tcg{(umelen-am)}	&	sun, day, sky (*j > {\0} in all \tsc{\ili{Ambon}-Seram})	\\
	&		&	*su\tbr{j}a	&	kana sau-a	&		&	spike trap (irr. *j > {\0})	\\
	&		&	*ulə\tbr{j}	&	ule	&	ula\tbr{n}-em	&	small worm (Liana-S. possibly \it{ul-e})	\\
\lspbottomrule
\end{tabularx}\mbox{}\\}

Similarly,
*l = \it{l} in most cases.
Two words have *l > \it{n} which may be conditioned by the adjacent high vowels (i.e.
*l > \it{n} /V\tsc{\tsc{[+high]}}{\gap}V\tsc{\tsc{[+high]}}).
These words are *kulit > **ulit-a > **ilut-a (vowel metathesis)
> **inut-a > liana-\ili{Seti} \it{int-a},
\ili{Benggoi} \it{int-am} ‘skin’ and *qulu > **unu-a > \it{un-a} ‘head’.
However, *l > \it{n} does not occur in other words with adjacent
high vowels *bulu > \ili{Liana-Seti} \it{fulaʔa},
\ili{Benggoi} \it{kafult-em} ‘body hair’,
or in *qulu > Adabai \ili{Seti} \it{ulu-,} \ili{Benggoi} \it{ulu-m} ‘head’.
\\

\noindent{\wg\st{6pt}\begin{tabularx}
{\linewidth}{>{\hspace{-\tabcolsep}}l<{\hspace{-\tabcolsep}}
>{\hspace{-\tabcolsep}\enspace}lll
>{\raggedright\arraybackslash}X}
\lsptoprule
%\tbx{10-63}	&	PMP	&	Liana-S.	&	\ili{Benggoi}	&	gloss	\\	\midrule
%\endfirsthead
\tbx{10-63}	&	PMP	&	Liana-S.	&	\ili{Benggoi}	&	gloss	\\	\midrule
%\endhead
	&	*\tbr{l}ayaR	&	\tbr{l}aiyal-a	&	\tbr{l}eh-am	&	sail (n.)	\\
	&	*\tbr{l}ima	&	en-\tbr{l}ima	&	en-\tbr{l}im	&	five	\\
	&	*\tbr{l}ahud	&	\tbr{l}au	&		&	sea, seaward	\\
	&	*\tbr{l}əsuŋ	&	\tbr{l}usn-a	&	kan\tbr{l}usn-am	&	mortar	\\
	&	*\tbr{l}iaŋ	&	\tbr{l}ia	&		&	cave	\\
	&	*za\tbr{l}an	&	la\tbr{l}a	&	la\tbr{l}ean	&	path	\\
	&	**zə\tbr{l}aq	&	le\tbr{l}a-k	&	le\tbr{l}a-m	&	tongue (PWMP)	\\
	&	*sa\tbr{l}aq	&	sa\tbr{l}a	&		&	wrong, fault	\\
	&	*bu\tbr{l}u	&	fu\tbr{l}aʔa (C),	&	kafu\tbr{l}t-em	&	body hair	\\	\hhline{~}
	&		&	fuwal-e	&		&	\hp{\,}	\\
	&	*u\tbr{l}əj	&	u\tbr{l}e	&	u\tbr{l}an-em	&	small worm	\\
	&	*hu\tbr{l}aR	&	u\tbr{l}a (C)	&	u\tbr{l}a-m	&	snake (\ili{Seti} Adabai dialect)	\\
	&	*tə\tbr{l}u	&	en-to\tbr{l}	&	en-to\tbr{l}	&	three	\\
	&	*wa\tbr{l}u	&	wa\tbr{l}	&	a\tbr{l}u	&	eight	\\
	&	**ma\tbr{l}ip	&	ma{\tl}ma\tbr{l}	&	mo\tbr{h}	&	laugh (\ili{Benggoi} possibly non-cognate)	\\
	&	**ma\tbr{l}ip	&	mum\tbr{l}iya	&		&	smile	\\
	&	*qatə\tbr{l}uR	&	to\tbr{ll}a	&	to\tbr{l}-am	&	egg	\\
	&	*ba\tbr{l}abaw 	&	u\tbr{l}afa	&		&	rat, mouse	\\
	&	*qu\tbr{l}u	&	u\tbr{l}u- (C)	&	u\tbr{l}u-m	&	head (\ili{Seti} Adabai dialect)	\\
	&	*qu\tbr{l}u	&	u\tbr{n}-a	&		&	head	\\
	&	*ku\tbr{l}it	&	i\tbr{n}t-a	&	i\tbr{n}t-am	&	skin (also Collins [1986] \it{inta})	\\
	&	**kawi\tbr{l}, **kabi\tbr{l}	&	ai, ai\tbr{l}-a	&		&	fish hook	\\
	&	*qa\tbr{l}ima	&	ima-k	&	ima-m	&	hand (irr. *l > {\0}, also in \ili{Manusela})	\\
\lspbottomrule
\end{tabularx}\mbox{}\\}

\newpage
In most cases *R > \it{l},
as shown in \qf{10-64},
with *R > \it{l,} {\0} word finally.\footnote{
		One instance of
		*R > {\0} word medially is \ili{Liambata} *duRi > \it{lu-im} ‘bone’
		\citep[168]{Stresemann1927}.}
\\

\noindent{\wg\st{4pt}\begin{tabularx}
{\linewidth}{>{\hspace{-\tabcolsep}}l<{\hspace{-\tabcolsep}}
>{\hspace{-\tabcolsep}\enspace}l<{\hspace{-\tabcolsep}\,}>{\hspace{-\tabcolsep}\,}lll
>{\raggedright\arraybackslash}X}
\lsptoprule
%\tbx{10-64}	&		&	PMP	&	Liana-S.	&	\ili{Benggoi}	&	gloss	\\	\midrule
%\endfirsthead
\tbx{10-64}	&		&	PMP	&	Liana-S.	&	\ili{Benggoi}	&	gloss	\\	\midrule
%\endhead
	&		&	*\tbr{R}usuk	&	\tbr{l}usu-k	&	\tbr{l}usu-m	&	rib	\\
	&		&	*qaba\tbr{R}a	&	maba\tbr{l}a-k (C),	&	mab\tbr{l}er	&	arm, shoulder	\\	\hhline{~}
	&		&		&	naba\tbr{l}a-k	&		&		\\
	&		&	*u\tbr{R}at	&	u\tbr{l}a tawal-e	&	\tcg{(urat)}	&	vein, tendon (\ili{Benggoi} from \ili{Malay})	\\
	&		&	*tu\tbr{R}un	&	au-ru\tbr{l}	&	\tcg{(dudu)}	&	I descend (onset alternation \srf{07-04-03})	\\
	&		&	*ta\tbr{R}aq	&	au-ra\tbr{l}a	&		&	I fell (a tree) (onset alternation \srf{07-04-03})	\\
	&		&	*baqə\tbr{R}u	&	fo\tbr{ll}-a	&	hol{\tl}ho\tbr{l}-am	&	new	\\
	&		&	*ma-bə\tbr{R}əqat	&	nefe\tbr{l}a,	&	he\tbr{l}an	&	heavy	\\	\hhline{~}
	&		&		&	mefe\tbr{l}a	&	&	\hp{\,}\\
	&		&	*pa\tbr{R}ih	&	fa\tbr{l}	&		&	stingray	\\	\midrule
\tbx{10-64b}	&	/{\gap}{\#}	&	*ndu\tbr{R}du\tbr{R} (?)	&	dó\tbr{l}-a (C)	&	\tcg{(bormuln-am)}	&	thunder	\\
	&		&	*timu\tbr{R}	&	tim\tbr{l}-a	&		&	\ili{east} monsoon	\\
	&		&	*qatəlu\tbr{R}	&	to\tbr{l}l-a	&	to\tbr{l}-am	&	egg	\\
	&		&		&	to\tbr{l}-n-im	&		&	egg (\ili{Liambata}, \citealp[22]{Stresemann1927})	\\
	&		&	*wahiyə\tbr{R}	&	wai-a	&	wei\tbr{l}-am	&	water	\\
	&		&	*busu\tbr{R}	&	usu-a	&	hus\tbr{l}-am	&	bow (*bu > u, also in \ili{Geser})	\\
	&		&	*hula\tbr{R}	&	ula (C)	&	ula-m	&	snake (\ili{Seti} Adabai)	\\
	&		&	*niu\tbr{R}	&	nu-a	&	\tcg{(hatey-am)}	&	coconut	\\
\lspbottomrule
\end{tabularx}\mbox{}\\}

\newpage
However, there are also words which have *R > \it{h},
{\0}, \it{ʔ} often with a different outcome in different varieties
of \tsc{Seti} and/or \ili{Benggoi}.
\citet[127f]{Collins1986} suggests that the reflex of *R may
be conditioned by the following vowel,
with *R > \it{l} before high vowels
and *R > \it{h} before non-high vowels.
However, there are also four examples of *R > \it{l} before high
vowels (*Rusuk, *duRi, *baqəRu, *paRih).
A few words, which are likely loans have, *R > \it{r}.
Examples in which *R does not become \it{l} are shown in \qf{10-65}.
\mbox{}\\

\noindent{\wg\st{2.5pt}\begin{tabularx}
{\linewidth}{>{\hspace{-\tabcolsep}}l<{\hspace{-\tabcolsep}}
>{\hspace{-\tabcolsep}\enspace}lll
>{\raggedright\arraybackslash}X}
\lsptoprule
%\tbx{10-65}	&	PMP	&	Liana-S.	&	\ili{Benggoi}	&		\\	\midrule
%\endfirsthead
\tbx{10-65}	&	PMP	&	Liana-S.	&	\ili{Benggoi}	&		\\	\midrule
%\endhead
	&	*da\tbr{R}aq	&	la\tbr{h}-a {\tl} la\tbr{h}-e	&	la\tbr{h}-im	&	blood	\\
	&	*ha\tbr{R}ezan	&	olam (C)	&	a\tbr{h}el-am	&	stairs (Liana-S. *R > {\0})	\\
	&	*wa\tbr{R}əj	&	ai we\tbr{h}el-a	&	iwel-am	&	vine, rope (\ili{Benggoi} *R > {\0}, *j > \it{l})	\\
	&	*ti\tbr{R}əm	&	te\tbr{h}en-a	&		&	oyster	\\
	&	*laya\tbr{R}	&	laiya\tbr{l}-a	&	le\tbr{h}-am	&	sail (n.)	\\
	&	**ŋa\tbr{R}ək	&	na\tbr{ʔ}a	&	\tcg{(tahun)}	&	year (cf. \ili{Buru} \it{ŋahe-t}, \ili{Asilulu} \it{nale} among others)	\\
	&	*ku\tbr{R}ita	&	u\tbr{w}ita ?/uita/	&		&	octopus (cf. \ili{Bobot} *R > w)	\\
	&	*ma\tbr{R}uqanay	&	mu\tbr{l}a(la)ina	&	mo\tbr{l}nan-am	&	male	\\
	&	*ma\tbr{R}uqanay	&	‹múana›	&		&	man (\ili{Liambata}, \citealp[617]{Wallace1869}	\\
	&	*du\tbr{R}i	&	lui\tbr{l}-e	&	lui\tbr{l}-am	&	bone	\\
	&	*du\tbr{R}i	&	lu-im ?/lui-m/	&		&	bone (\ili{Liambata}, \citealp[135]{Stresemann1927}	\\
	&	*\tbr{R}umbia	&	ipiya	&	(bam)	&	sago (*R > {\0} in many languages)	\\
	&	*haba\tbr{R}at	&	fa\tbr{r}at-a	&	\tcg{(barat)}	&	Liana-S. `west monsoon', \ili{Benggoi} `west' < \ili{Malay}	(irr. *R > \it{r}, also in \tsc{East Littoral} \srf{10-01-02}	\\
	&	*qa\tbr{R}us	&	a\tbr{r}us	&	a\tbr{r}us	&	current (in water) (irr. *R > \it{r}, \ili{Malay} \it{arus}?)	\\
\lspbottomrule
\end{tabularx}\mbox{}\\}

\newpage
There are a number of reflexes of PMP reconstructions with *R
in the ``\ili{Liambata}'' data in \citet{Tauern1931} which are either not
attested in other sources or which show differences to other sources.
All reflexes of reconstructions with *R in \citet{Tauern1931}
which have been identified are given in \qf{10-66}.
Note in particular the reflexes of *wahiyəR ‘water’
and **kambeR which have *R > \it{h}.\footnote{
		For **kambeR ‘saliva’ compare Proto-Rote-\ili{Meto}
		*kambe ‘saliva’ \citep[192]{Edwards2021},
		\ili{Kamarian} \it{aper} ‘mucus’, \ili{Nusalaut} \it{apel} ‘mucus,
		snot’, and South \ili{Nuaulu} \it{apai},
		\it{apana} ‘saliva’.}
\\

\noindent{\wg\st{6pt}\begin{tabularx}
{\linewidth}{>{\hspace{-\tabcolsep}}l<{\hspace{-\tabcolsep}}
>{\hspace{-\tabcolsep}\enspace}ll
>{\raggedright\arraybackslash}X}
\lsptoprule
%\tbx{10-66}	&	PMP	&	\ili{Liambata}	&	gloss	\\	\midrule
%\endfirsthead
\tbx{10-66}	&	PMP	&	``\ili{Liambata}''	&	gloss	\\	\midrule
%\endhead
	&	*qaba\tbr{R}a	&	fa\tbr{l}a-n	&	hand	\\
	&	*bə\tbr{R}əqat	&	fe\tbr{l}a-n	&	heavy	\\
	&	*qatəlu\tbr{R}	&	to\tbr{l}-n-a	&	egg	\\
	&	*\tbr{R}ibun	&	\tbr{l}ifin=sa	&	one thousand	\\
	&	*du\tbr{R}i	&	ian li\tbr{l}i-na	&	fish bone	\\
	&	*wahiyə\tbr{R}	&	we\tbr{h}en-a	&	water, broth (possibly \it{weh-en-a})	\\
	&	*wahiyə\tbr{R}	&	wei\tbr{l}-a, wai\tbr{l}-a	&	water	\\
	&	*laya\tbr{R}	&	lilaian-a	&	sail (cf. \ili{Liana-Seti} \it{laiyal-a})	\\
	&	**kambe\tbr{R}	&	epe\tbr{h}e	&	saliva (possibly \it{epeh-e})	\\
	&	*\tbr{R}umbia	&	bia	&	sago	\\
	&	*liqə\tbr{R}	&	lia-n	&	voice	\\
\lspbottomrule
\end{tabularx}\mbox{}\\}

The reflexes of schwa are relatively complex.
Penultimate schwa is usually reflected as \it{o},
though there are also instances of penultimate *ə > \it{e},
often in words which also have unexpected *ə > \it{e} in other \tsc{\ili{Ambon}-Seram} languages.
Schwa in final syllables shows a split in \tsc{Seti} with the most
common reflex being \it{e} or \it{a}.
Most cases of *ə > \it{e} in final syllables are after penultimate
\it{e} and may be due to vowel harmony.
Note that *ə > \it{a} in final syllables is characteristic of
neighbouring \tsc{Seram-Tanimbar-Bomberai} languages (\srf{11-01-02}),
though this change is much more regular in these \tsc{Seram-Tanimbar-Bomberai}
languages than it is in \tsc{Seti}.

Final schwa is usually lost in \ili{Benggoi},
with only one putative retention identified: *quləj > \it{ul\tbr{a}n-am}
‘worm’ which also shows unexpected *j > \it{n}.
Examples of \ili{Benggoi} *ə > {\0} in final syllables not shown in
\qf{10-69} include *bəqbəq > \it{hoh-am} ‘mouth’
and *təkən > \it{ai ton-am} ‘staff’.

\noindent{\wg\st{3pt}\begin{tabularx}
{\linewidth}{>{\hspace{-\tabcolsep}}l<{\hspace{-\tabcolsep}}
>{\hspace{-\tabcolsep}\enspace}l<{\hspace{-\tabcolsep}\,}>{\hspace{-\tabcolsep}\,}lll
>{\raggedright\arraybackslash}X}
\lsptoprule
%\tbx{10-67}	&		&	PMP	&	Liana-S.	&	\ili{Benggoi}	&	gloss	\\	\midrule
%\endfirsthead
\tbx{10-67}	&		&	PMP	&	Liana-S.	&	\ili{Benggoi}	&	gloss	\\	\midrule
%\endhead
	&	*ə > \it{o}	&	*qat\tbr{ə}luR	&	t\tbr{o}lla	&	t\tbr{o}l-am	&	egg	\\
	&		&	*t\tbr{ə}lu	&	en-t\tbr{o}l	&	en-t\tbr{o}l	&	three	\\
	&		&	*b\tbr{ə}taw	&	f\tbr{o}to-nin	&	b\tbr{o}t-am	&	opposite sex sibling (irr. *-aw > \it{o})	\\
	&		&	*baq\tbr{ə}Ru	&	f\tbr{o}lla	&	hol{\tl}h\tbr{o}l-am	&	new	\\
	&		&	*t\tbr{ə}buh	&	t\tbr{o}f-a	&	\tcg{(lakl-am)}	&	sugarcane	\\
	&		&	*ma-p\tbr{ə}nuq	&	b\tbr{o}nu (C)	&		&	full	\\
	&		&	**t\tbr{ə}pu (PCM)	&	t\tbr{o}h-a	&	\tcg{(ato)}	&	chicken	\\
	&		&	**mb\tbr{ə}tu	&	b\tbr{o}tu-a	&	\tcg{(noni)}	&	night	\\
	&		&	*k\tbr{ə}dəŋ	&	au k-\tbr{o}lo	&		&	stand	\\
	&		&	*haR\tbr{ə}zan	&	\tbr{o}lam (C)	&	ah\tbr{e}l-am	&	stairway	\\	\midrule
\tbx{10-68}	&	*ə > \it{e}	&	*qal\tbr{ə}jaw	&	l\tbr{e}a	&	\tcg{(umelen-am)}	&	sun, day	\\
	&		&	*p\tbr{ə}ñu	&	f\tbr{e}n-a, f\tbr{eⁱ}n-a	&		&	turtle	\\
	&		&	**l\tbr{ə}si (CMP)	&	l\tbr{e}sn-e	&		&	excess	\\
	&		&	*q\tbr{ə}nay	&	\tbr{e}nn-e	&	\tbr{i}nen-am	&	sand	\\
	&		&	*ma-b\tbr{ə}Rəqat	&	nef\tbr{e}la,	&	h\tbr{e}lan	&	heavy	\\	\hhline{~}
	&		&	&	mef\tbr{e}la	&		&	\hp{\,}	\\	\midrule
\tbx{10-69}	&	*ə /σ\und{σ}	&	*ul\tbr{ə}j	&	ul\tbr{e} 	&	ul\tbr{a}n-em	&	small worm  (possibly \it{ul-e})	\\
	&		&	*waR\tbr{ə}j	&	ai weh\tbr{e}l-a	&	iwel-am	&	vine, rope	\\
	&		&	*ma-qit\tbr{ə}m	&	met\tbr{e}n-a	&	\tcg{(lobot)}	&	black (*ə > \it{e} in all \tsc{\ili{Ambon}-Seram})	\\
	&		&	*tiR\tbr{ə}m	&	teh\tbr{e}n-a (C) 	&		&	oyster	\\
	&		&	*tuq\tbr{ə}d	&	tu\tbr{a}-ne	&		&	tree stump	\\
	&		&	*qin\tbr{ə}p	&	au k-in\tbr{a}	&	\tcg{(let)}	&	I sleep	\\
	&		&	*qat\tbr{ə}p	&	at\tbr{a}-lanne	&	etlan-am	&	thatch, roof	\\
	&		&	*kəd\tbr{ə}ŋ	&	au k-ol\tbr{o}	&		&	stand (vowel harmony?)	\\
	&		&	*dak\tbr{ə}p	&	au-lak	&		&	I catch, seize	\\
\lspbottomrule
\end{tabularx}\mbox{}\\}

\newpage
\tsc{Seti} appears to have a split of *ma-p/*mb > \it{b,
p} and \ili{Benggoi} has *ma-p/*mb > \it{b}
in the two identified reflexes.
Examples as shown in \qf{10-70}.
This is different to other \tsc{\ili{Ambon}-Seram} languages which
consistently have a devoiced reflex; \it{mp,} and/or \it{p}.
\\

\noindent{\wg\st{6pt}\begin{tabularx}
{\linewidth}{>{\hspace{-\tabcolsep}}l<{\hspace{-\tabcolsep}}
>{\hspace{-\tabcolsep}\enspace}lll
>{\raggedright\arraybackslash}X}
\lsptoprule
%\tbx{10-70}	&	PMP	&	Liana-S.	&	\ili{Benggoi}	&	gloss	\\	\midrule
%\endfirsthead
\tbx{10-70}	&	PMP	&	Liana-S.	&	\ili{Benggoi}	&	gloss	\\	\midrule
%\endhead
	&	*\tbr{ma-p}ənuq	&	\tbr{b}onu (C)	&		&	full	\\
	&	**\tbr{mb}ətu		&	\tbr{b}otu-a		&	\tcg{(noni)}	&	night	\\
	&	*\tbr{ma-p}utiq	&	\tbr{b}utt-a		&	but{\tl}\tbr{b}ut-am	&	white	\\
	&	*\tbr{ma-p}utiq	&	\tbr{p}ut-a (C)	&		&	white	\\
	&	*Ru\tbr{mb}ia		&	i\tbr{p}iya			&	\tcg{(bam)}	&	sago tree (cf. ``\ili{Liambata}'' \it{bia} `sago' [\citealp[132]{Tauern1931}])	\\
	&	*ka\tbr{mb}u		&		&	a\tbr{b}un	&	pregnant	\\
\lspbottomrule
\end{tabularx}\mbox{}\\}

To summarise,
\tsc{Seti} and \ili{Benggoi} group with \tsc{\ili{Ambon}-Seram} on the basis of
*ŋ/*n > \it{n} and the usual merger of *d/*z/*l/*j > \it{l.} However,
the reflexes of *R do not fully merge with *d/*z/*l/*j
and there is a split before non-high vowels of *R > \it{l, h}.
\tsc{Seti} also shows voiced reflexes of *ma-p,
which is different to most \tsc{\ili{Ambon}-Seram}.
We place \tsc{Seti} and \ili{Benggoi} together as a primary branch of \tsc{\ili{Ambon}-Seram}
and agree with \citet[135]{Collins1986} that more data on varieties
of \tsc{Seti} is imperative.
